\section{Fundamentação}

\textbf{Máquinas de vetores de suporte.} A SVM~\cite{boser:92,cortes:95} separa duas classes por um hiperplano $f(x) = w \cdot x + b$, respondendo a classe positiva quando $f(x) > 0$. Entre os hiperplanos que separam o treino, ela escolhe o de maior margem, a faixa vazia entre a fronteira e os pontos mais próximos de cada classe, de largura $2/\lVert w \rVert$. Como dados reais raramente se separam por completo, a formulação de margem suave aceita pontos dentro da margem ou do lado errado, cobrando por eles a perda \emph{hinge}:
\begin{equation*}
  \min_{w,\,b} \; \tfrac{1}{2}\lVert w \rVert^2 + C \sum_i \max\bigl(0,\; 1 - y_i\, f(x_i)\bigr), \qquad y_i \in \{-1, +1\}.
\end{equation*}
O hiperparâmetro $C$ pesa os erros contra a largura da margem: $C$ pequeno tolera erros e deixa a fronteira simples; $C$ grande tenta acertar todo o treino. A solução depende só dos pontos com perda positiva ou sobre a margem, os \emph{vetores de suporte}. A SVM não estima probabilidades: sua saída é a distância com sinal até a fronteira, que serve para ordenar exemplos e calcular a AUC. Probabilidades exigem uma calibração à parte~\cite{platt:99}.

\textbf{Kernels.} No problema dual, os dados aparecem só em produtos $x_i \cdot x_j$. Trocar esse produto por um kernel $K(x_i, x_j)$ equivale a traçar o hiperplano num espaço transformado, sem calculá-lo, e a fronteira no espaço original fica curva: $f(x) = \sum_i \alpha_i y_i K(x_i, x) + b$, com a soma sobre os vetores de suporte. Usamos o kernel polinomial, $(\gamma\, x \cdot x' + c_0)^d$, e o RBF, $\exp(-\gamma \lVert x - x' \rVert^2)$. No RBF, $\gamma$ controla o alcance de cada vetor de suporte: pequeno, a fronteira é suave; grande, ela contorna os pontos do treino. Como o RBF mede distâncias, atributos em escalas diferentes precisam ser padronizados~\cite{hsu:03}. O treino pelo LIBSVM~\cite{chang:11} trabalha com pares de exemplos, e seu custo cresce entre o quadrado e o cubo do tamanho do treino; a previsão compara cada exemplo novo com todos os vetores de suporte. Para o kernel linear, o LIBLINEAR~\cite{fan:08} otimiza diretamente os pesos $w$, com custo proporcional ao tamanho do treino.

\textbf{Redes neurais.} Um neurônio artificial calcula uma soma ponderada das entradas, mais um viés, e aplica a ela uma função de ativação. Com a sigmoide, que leva qualquer número a um valor entre 0 e 1, um único neurônio é a regressão logística. O perceptron de múltiplas camadas (MLP) empilha camadas de neurônios entre a entrada e a saída, as camadas ocultas; com a ativação ReLU, $\max(0, z)$~\cite{nair:10}, elas permitem fronteiras de decisão não lineares. Os pesos começam sorteados e são ajustados pela retropropagação do erro~\cite{rumelhart:86}: a cada lote de exemplos, calcula-se quanto cada peso contribuiu para a perda (aqui, a entropia cruzada) e move-se o peso no sentido de reduzi-la, com um passo dado pela taxa de aprendizado. O otimizador Adam~\cite{kingma:15} adapta esse passo para cada peso. Uma época é uma passada completa pelo treino. Como o passo é o mesmo para entradas de escalas muito diferentes, os atributos costumam ser padronizados antes~\cite{lecun:98}. Treinar por épocas demais leva a rede a se ajustar a detalhes do treino; a parada antecipada~\cite{prechelt:98} separa uma parte do treino como validação e interrompe o treino quando o desempenho nela para de melhorar. Ao contrário da SVM, a rede é treinada para estimar probabilidades. Em dados tabulares, modelos de árvore ainda costumam superar redes neurais~\cite{grinsztajn:22}.

\textbf{Modelos de referência.} O CART~\cite{breiman:84} divide o espaço de atributos de forma recursiva, em cada nó escolhendo o limiar de um atributo que mais reduz a impureza (Gini ou entropia); a vantagem é a legibilidade das regras. Uma árvore de profundidade 1 é uma \emph{regra de um corte}: um único atributo comparado com um único limiar. A Random Forest~\cite{breiman:01} treina muitas árvores, cada uma numa amostra com reposição do treino (\emph{bootstrap}) e com sorteio dos atributos de cada nó, e a previsão é a média dos votos. A regressão logística, como visto acima, equivale a uma rede de um único neurônio.
